\subsection{对称希尔伯特幂}

设 E 为希尔伯特空间，n 为正整数。以 \(\hat{\mathbf{T}}^{n}(\mathrm{E})\) 或 \(E^{\otimes_{n}}\) 表示 n 个都等于 E 的希尔伯特空间的张量积。换言之，\(\hat{\mathbf{T}}^{n}(\mathrm{E})\) 是空间 \(\mathbf{T}^{n}(\mathrm{E}) = \mathrm{E} \otimes \ldots \otimes \mathrm{E}\)（n 个因子）关于由下式定义的分离准希尔伯特空间结构的完备化

\[
\left\langle x _ {1} \otimes \dots \otimes x _ {n} \mid y _ {1} \otimes \dots \otimes y _ {n} \right\rangle = \prod_ {i = 1} ^ {n} \left\langle x _ {i} \mid y _ {i} \right\rangle .
\]

若 \((e_i)_{i\in I}\) 是 \(E\) 的标准正交基，则向量族 \(e_{i_1} \otimes \ldots \otimes e_{i_n}\)（\(i_1, \ldots, i_n\) 取遍 \(I\)）是 \(\hat{\mathbf{T}}^n(E)\) 的标准正交基（V, p.~29, 推论 1）。我们得到 \(\hat{\mathbf{T}}^\circ(E) = K\)。

设 \(\sigma \in \mathfrak{S}_n\) 为集合 \(\{1, 2, ..., n\}\) 的一个置换。由 V, p.~27，存在自同构 \(p_{\sigma}\)（\(\hat{\mathbf{T}}^n(\mathbf{E})\) 的），它由下式刻画

\[
p _ {\sigma} (x _ {1} \otimes \dots \otimes x _ {n}) = x _ {\sigma^ {- 1} (1)} \otimes \dots \otimes x _ {\sigma^ {- 1} (n)}.
\]

我们有 \(p_{\sigma \tau} = p_{\sigma}p_{\tau}\)（\(\sigma, \tau\) 属于 \(\mathfrak{S}_n\)），因而自同态 \(\Pi_n = \frac{1}{n!}\sum_{\sigma \in \mathfrak{S}_n}p_\sigma\)（向量空间 \(\hat{\mathbf{T}}^n (\mathrm{E})\) 的）是到在 \(\mathfrak{S}_n\) 下保持不变的所有元素组成的子空间上的正交投影算子。然而（A, III, § 6, No.~3），\(\Pi_n\) 把« 代数 »张量积 \(\mathbf{T}^n (\mathrm{E})\) 映到子空间 \(\mathbf{T}\mathbf{S}^n (\mathrm{E})\)（由所有 \(n\) 阶对称张量组成的）上。换言之，\(\Pi_n\) 的像是空间 \(\mathbf{T}\mathbf{S}^n (\mathrm{E})\)（赋以由 \(\mathbf{T}^n (\mathrm{E})\) 的内积诱导的内积）的完备化；这一完备化将记作 \(\widehat{\mathbf{T}\mathbf{S}}^n (\mathrm{E})\)。

设 \(\mathbf{S}^n (\mathrm{E})\) 为向量空间 E 的 \(n\) 次对称幂（A, III, § 6, No.~1）。从 \(\mathbf{T}^n (\mathrm{E})\) 到 \(\mathbf{S}^n (\mathrm{E})\) 上的典范映射通过限制定义了一个同构 \(\lambda_{n}\)（从 \(\mathbf{TS}^n (\mathrm{E})\) 到 \(\mathbf{S}^n (\mathrm{E})\) 上）。我们立即可以验证，其逆同构由下式给出

\[
\mu_ {n} (x _ {1} \dots x _ {n}) = \Pi_ {n} (x _ {1} \otimes \dots \otimes x _ {n}) = \frac {1}{n !} \sum_ {\sigma \in \mathfrak {S} _ {n}} x _ {\sigma^ {- 1} (1)} \otimes \dots \otimes x _ {\sigma^ {- 1} (n)}\tag{13}
\]

其中 \(x_{1},\ldots ,x_{n}\) 属于 E。

我们通过令下式成立而在 \(\mathbf{S}^{n}(\mathrm{E})\) 上定义一个分离准希尔伯特空间结构

\[
\langle u | v \rangle = n! \left\langle \mu_ {n} (u) | \mu_ {n} (v) \right\rangle .\tag{14}
\]

于是我们有（试与 A, III, § 11, No.~5 的公式 (29) 比较）

\[
\left\langle x _ {1} \dots x _ {n} \mid y _ {1} \dots y _ {n} \right\rangle = \sum_ {\sigma \in \mathfrak {S} _ {n}} \prod_ {i = 1} ^ {n} \left\langle x _ {i} \mid y _ {\sigma (i)} \right\rangle ,\tag{15}
\]

特别地

\[
\langle x ^ {n} | y ^ {n} \rangle = n! \langle x | y \rangle^ {n}.\tag{16}
\]

设 \(\hat{\mathbf{S}}^n (\mathrm{E})\) 表示准希尔伯特空间 \(\mathbf{S}^n (\mathrm{E})\) 的完备化，\(\hat{\mathbf{S}} (\mathrm{E})\) 表示诸希尔伯特空间 \(\hat{\mathbf{S}}^n (\mathrm{E})\) 的外希尔伯特和。我们可以证明（V, p.~73, 练习 1），除非 E 恰为 0，代数 \(\mathbf{S}(\mathrm{E})\) 中的乘法不可能连续地延拓到 \(\mathbf{S}(\mathrm{E})\) 上。

命题 4. --- 设 \((e_{i})_{i\in I}\) 为希尔伯特空间 E 的标准正交基。对每个 \(\alpha\)（属于 \(\mathbf{N}^{(1)}\)），令

\[
z _ {\alpha} = \prod_ {i \in I} e _ {i} ^ {\alpha_ {i}} / (\alpha_ {i}!) ^ {1 / 2}.\tag{17}
\]

于是 \((z_{\alpha})_{\alpha \in \mathbf{N}^{(1)}}\) 是 \(\hat{\mathbf{S}} (\mathrm{E})\) 的标准正交基。

设 \(E_{0}\) 为 E 中由向量 \(e_{i}\)（i 取遍 I）生成的向量子空间。于是诸 \(z_{\alpha}\) 构成向量空间 \(\mathbf{S}(E_{0})\) 的基（A, III, § 6, No.~6）。但 \(E_{0}\) 在 E 中稠密，且多线性映射 \((x_{1},...,x_{n})\mapsto x_{1}\ldots x_{n}\)（从 \(E\times\cdots\times E\) 到 \(\mathbf{S}(E)\) 中）对一切 \(n\geqslant1\) 是连续的；因此 \(\mathbf{S}(E_{0})\) 在 \(\mathbf{S}(E)\) 中稠密。现在只需证明诸 \(z_{\alpha}\) 组成的族是标准正交的。首先注意，\(\hat{\mathbf{S}}^{n}(\mathrm{E})\) 与 \(\hat{\mathbf{S}}^{m}(\mathrm{E})\)（当 \(n \neq m\) 时）正交。于是只需证明公式

\[
\langle z _ {\alpha} | z _ {\beta} \rangle = \left\{ \begin{array}{l l} 1 & \text {if} \quad \alpha = \beta \\ 0 & \text {if} \quad \alpha \neq \beta \end{array} \right.
\]

其中 \(|\alpha| = \sum_{i \in I} \alpha_i\) 与 \(|\beta| = \sum_{i \in I} \beta_i\) 等于同一个整数 \(n\)。

考虑一个划分 \((\mathbf{P}_i)_{i\in I}\)（集合 \(\{1,2,\dots,n\}\) 的划分），使得 \(\operatorname{Card} \mathbf{P}_i = \alpha_i\) 对所有 \(i\in I\) 成立。令 \(x_k = e_i\)（若 \(k\) 属于 \(\mathbf{P}_i\)），于是 \(x_1\ldots x_n = \prod_{i\in I}e_i^{\alpha_i}\)。类似地，我们定义 \((\mathbf{Q}_i)_{i\in I}\) 与 \(y_k\)，使得 \(\operatorname{Card} \mathbf{Q}_i = \beta_i\) 且 \(y_1\ldots y_n = \prod_{i\in I}e_i^{\beta_i}\)。由于诸 \(e_i\) 两两正交，我们有 \(\langle x_k|y_{\sigma(k)}\rangle = 0\)，除非存在指标 \(i\in I\) 使 \(k\in \mathbf{P}_i\) 且 \(\sigma (k)\in \mathbf{Q}_i\)。于是由公式 (15)，我们有 \(\langle x_1\ldots x_n|y_1\ldots y_n\rangle = 0\)，除非存在置换 \(\sigma \in \mathfrak{S}_n\) 使 \(\sigma (\mathbf{P}_i) = \mathbf{Q}_i\) 对所有 \(i\in I\) 成立（这蕴含 \(\alpha = \beta\)）。于是 \(\langle z_{\alpha}|z_{\beta}\rangle = 0\)（当 \(\alpha \neq \beta\) 时）。同样的论证证明 \(\| x_1\ldots x_n\|^2\) 等于 \(\sigma \in \mathfrak{S}_n\) 中使 \(\sigma (\mathbf{P}_i) = \mathbf{P}_i\) 对所有 \(i\in I\) 成立者的个数，从而等于 \(\prod_{i=1}^{\infty}\alpha_i!\)。由此得 \(\| z_\alpha \| = 1\)，命题得证。

推论. --- 假定希尔伯特空间 E 是正交子空间 M 与 N 的直和。从 \(\mathbf{S}(\mathbf{M}) \otimes \mathbf{S}(\mathbf{N})\) 到 \(\mathbf{S}(\dot{\mathbf{E}})\) 上的典范同构 g（A, III, §6, No.~6）可唯一地延拓为从 \(\hat{\mathbf{S}}(\mathbf{M}) \otimes_{2} \hat{\mathbf{S}}(\mathbf{N})\) 到 \(\hat{\mathbf{S}}(\mathbf{E})\) 上的希尔伯特空间同构 h。

设 \((e_{i})_{i\in I}\)（相应地 \((f_{j})_{j\in J}\)）为希尔伯特空间 M（相应地 N）的标准正交基，并设 \(M_{0}\)（相应地 \(N_{0}\)）为 E 中由向量 \(e_{i}\)（相应地 \(f_{j}\)）生成的向量子空间。令 \(E_{0}=M_{0}+N_{0}\)，并以 \(g_{0}\) 表示从 \(S(M_{0})\otimes S(N_{0})\) 到 \(S(E_{0})\) 上的典范同构。令

\[
z _ {\alpha} = \prod_ {i \in \mathbf {I}} e _ {i} ^ {\alpha_ {i}} / (\alpha_ {i}!) ^ {1 / 2}, t _ {\beta} = \prod_ {j \in \mathbf {J}} f _ {j} ^ {\beta _ {j}} / (\beta_ {j}!) ^ {1 / 2}
\]

其中 \(\alpha \in \mathbf{N}^{(I)}\)，\(\beta \in \mathbf{N}^{(J)}\)。由命题 4，我们这样就定义了标准正交基 \((z_{\alpha})_{\alpha \in \mathbf{N}^{(I)}}\)（\(\hat{\mathbf{S}}(\mathbf{M})\) 的）、\((t_{\beta})_{\beta \in \mathbf{N}^{(J)}}\)（\(\hat{\mathbf{S}}(\mathbf{N})\) 的）以及 \((z_{\alpha}t_{\beta})_{\alpha \in \mathbf{N}^{(I)},\beta \in \mathbf{N}^{(J)}}\)（\(\hat{\mathbf{S}}(\mathbf{E})\) 的）。由于 \(z_{\alpha}t_{\beta} = g_0(z_{\alpha}\otimes t_{\beta})\)，且诸元素 \(z_{\alpha}\otimes t_{\beta}\) 构成 \(\hat{\mathbf{S}}(\mathbf{M})\hat{\otimes}_2\hat{\mathbf{S}}(\mathbf{N})\) 的标准正交基（V, p.~29, 推论 1），可见 \(g_0\) 可延拓为希尔伯特空间同构 \(h:\hat{\mathbf{S}}(\mathbf{M})\hat{\otimes}_2\hat{\mathbf{S}}(\mathbf{N})\to \hat{\mathbf{S}}(\mathbf{E})\)。根据构造，我们有

\[
h (x _ {1} \dots x _ {m} \otimes y _ {1} \dots y _ {n}) = x _ {1} \dots x _ {m} y _ {1} \dots y _ {n}
\]

其中向量 \(x_{1}, \ldots, x_{m}\) 属于 \(M_{0}\)，向量 \(y_{1}, \ldots, y_{n}\) 属于 \(N_{0}\)。由连续性，同一关系式对 M 中的向量 \(x_{1}, \ldots, x_{m}\) 与 N 中的向量 \(y_{1}, \ldots, y_{n}\) 也成立；换言之，h 延拓了 g。h 的唯一性是显然的。

设 E 与 F 为两个希尔伯特空间，\(u \in \mathcal{L}(E; F)\)。线性映射 \(\hat{\mathbf{T}}^{n}(u) = u \otimes_{2} \ldots \otimes_{2} u\)（n 个因子）是从 \(\hat{\mathbf{T}}^{n}(E)\) 到 \(\hat{\mathbf{T}}^{n}(F)\) 中的连续映射，其范数为 \(\|u\|^{n}\)（V, p.~28, 公式 (10)）。此外，V, p.~30 的公式 (13) 与 (14) 表明，存在一个同构 \(\phi_{n,\mathrm{E}}\)（从 \(\hat{\mathbf{S}}^n (\mathrm{E})\) 到子空间 \(\widehat{\mathbf{T}\mathbf{S}}^n (\mathrm{E})\)（\(\hat{\mathbf{T}}^n (\mathrm{E})\) 的）上）且仅存在一个，使得

\[
\phi_ {n, \mathrm{E}} (x _ {1} \dots x _ {n}) = \frac {1}{(n !) ^ {1 / 2}} \sum_ {\sigma \in \mathfrak {S} _ {n}} x _ {\sigma (1)} \otimes \dots \otimes x _ {\sigma (n)} \quad (x _ {1},..., x _ {n} \text { in } \mathrm{E}).\tag{18}
\]

因此存在一个连续线性映射 \(\hat{\mathbf{S}}^{n}(u)\)（从 \(\hat{\mathbf{S}}^{n}(\mathrm{E})\) 到 \(\hat{\mathbf{S}}^{n}(\mathrm{F})\) 中）且仅存在一个，它使下面的图表交换

\[
\begin{array}{c c c} \hat {\mathbf {S}} ^ {n} (\mathrm{E}) & \xrightarrow {\phi_ {n , \mathrm{E}}} & \hat {\mathbf {T}} ^ {n} (\mathrm{E}) \\ \hat {\mathbf {S}} ^ {n} (u) \Big \downarrow & & \Big \downarrow \hat {\mathbf {T}} ^ {n} (u) \\ \hat {\mathbf {S}} ^ {n} (\mathrm{F}) & \xrightarrow {\phi_ {n , \mathrm{F}}} & \hat {\mathbf {T}} ^ {n} (\mathrm{F}) \end{array}
\]

我们现在证明公式

\[
\left\| \hat {\mathbf {S}} ^ {n} (u) \right\| = \| u \| ^ {n}.\tag{19}
\]

显然有 \(\| \hat{\mathbf{S}}^n (u)\| \leqslant \| \hat{\mathbf{T}}^n (u)\| = \| u\| ^n\)。另一方面，对一切 \(x\in \mathrm{E}\)，有 \(\hat{\mathbf{S}}^n (u)(x^n) = u(x)^n,\| x^n\| = (n!)^{1 / 2}\| x\| ^n\) 以及 \(\| u(x)^n\| = (n!)^{1 / 2}\| u(x)\| ^n\)，由此得

\[
\left\| \hat {\mathbf {S}} ^ {n} (u) \right\| \| x \| ^ {n} \geqslant \left\| u (x) \right\| ^ {n};
\]

由此立即得出 \(\| \hat{\mathbf{S}}^n (u)\| \geqslant \| u\| ^n\)，因而公式 (19) 成立。

显然我们有公式

\[
\hat {\mathbf {S}} ^ {n} (1 _ {\mathrm{E}}) = 1 _ {\hat {\mathbf {S}} ^ {n} (\mathrm{E})}\tag{20}
\]

\[
\hat {\mathbf {S}} ^ {n} (v \circ u) = \hat {\mathbf {S}} ^ {n} (v) \circ \hat {\mathbf {S}} ^ {n} (u) \quad \text { for } \quad v \in \mathscr {L} (\mathrm{F}; \mathrm{G}).\tag{21}
\]

最后，\(\hat{\mathbf{S}}^n (u)\) 在 \(\mathbf{S}^n (\mathrm{E})\) 上与 A, III, § 6, No.~2 中定义的线性映射 \(\mathbf{S}^n (u):\mathbf{S}^n (\mathrm{E})\to \mathbf{S}^n (\mathrm{F})\) 一致，因为它把 \(x_{1}\ldots x_{n}\) 变为 \(u(x_{1})\ldots u(x_{n})\)（对 E 中的每个 \(x_{1},\ldots ,x_{n}\)）。

例. --- * 1) 设 \(d \geqslant 1\) 为整数，\(\omega\) 为 \(\mathbf{R}^{d}\) 上关于 Lebesgue 测度 \(\mu\) 局部可积的正函数。设 E 为希尔伯特空间 \(\mathrm{L}^{2}(\mathbf{R}^{d}, \omega, \mu)\)，并令 \(\mathbf{S} = \mathbf{S}(\mathrm{E})\)。于是 \(\mathbf{S}\) 可与所有序列 \(f = (f_{n})_{n \geqslant 0}\) 组成的空间等同，其中每个 \(f_{n}\) 是 \((\mathbf{R}^{d})^{n}\) 上的函数，它关于 Lebesgue 测度 \(\mu \otimes \ldots \otimes \mu\)（\(n\) 个因子）可测，且在 \(n\) 个因子（\((\mathbf{R}^{d})^{n}\) 的）的置换下不变，并且满足

\[
\| \mathbf {f} \| ^ {2} = \sum_ {n = 0} ^ {\infty} n! \int_ {\mathbf {R} ^ {d}} \dots \int_ {\mathbf {R} ^ {d}} \left| f _ {n} (\mathbf {x} _ {1}, \dots , \mathbf {x} _ {n}) \right| ^ {2} \omega (\mathbf {x} _ {1}) \dots \omega (\mathbf {x} _ {n}) d \mathbf {x} _ {1} \dots d \mathbf {x} _ {n}\tag{22}
\]

为有限。范数 \(\| \mathbf{f}\|\)（\(\mathbf{S}\) 中的）由公式 (22) 定义。上面定义的希尔伯特空间 \(\mathbf{S}\) 称为对应于权 \(\omega_{*}\) 的对称 Fock 空间

* 2) 设 X 为分离的拓扑空间，\(\mu\) 为 X 上范数为 1 的正测度，E 为实希尔伯特空间 \(L_{\mathbb{R}}^{2}(X,\mu)\) 的希尔伯特子空间。若下列等价条件之一满足，则称 E 为高斯空间：

\begin{enumerate}
\def\labelenumi{\alph{enumi})}
\item
  对所有 \(f \in \mathrm{E}\)，有 \(\int_{\mathbf{x}} e^{if} d\mu = \exp(-\|f\|^2/2)\)；
\item
  对所有范数为 1 的 \(f \in E\)，测度 \(\mu\) 在 f 下的像是测度
\end{enumerate}

\[
(2 \pi) ^ {- 1 / 2} e ^ {- x ^ {2} / 2} d x.
\]

在 \(\mathbf{R}\) 上。

假定 E 是高斯空间。设 \(f_1, \ldots, f_n\) 为若干函数，其类 \(f_i\) 属于 E。我们用下式在 X 上定义一个函数 :\(f_1 \ldots f_n\):（称为 \(f_1, \ldots, f_n\) 的« Wick 乘积 »）

\[
: f _ {1} \dots f _ {n} := \sum_ {0 \leqslant 2 p \leqslant n} (- 1) ^ {p} \sum_ {\sigma \in I _ {p}} \prod_ {i = 1} ^ {p} \left\langle f _ {\sigma (2 i - 1)} \mid f _ {\sigma (2 i)} \right\rangle \prod_ {j = 2 p + 1} ^ {n} f _ {\sigma (j)},\tag{23}
\]

其中 \(\mathbf{I}_p\) 是由满足下列条件的置换 \(\sigma\)（集合 \(\{1,2,\dots,n\}\) 的）组成的集合

\[
\begin{array}{l} \sigma (1) <   \sigma (2), \dots , \sigma (2 p - 1) <   \sigma (2 p) \\ \sigma (1) <   \sigma (3) <   \dots <   \sigma (2 p - 1) \\ \sigma (2 p + 1) <   \sigma (2 p + 2) <   \dots <   \sigma (n). \end{array}
\]

于是存在一个同构 \(\phi\)（从 \(\hat{\mathbf{S}}(\mathrm{E})\) 到 \(\mathrm{L}_{\mathbb{R}}^{2}(\mathrm{X},\mu)\) 的一个希尔伯特子空间上），它把乘积 \(\dot{f}_1\dots \dot{f}_n\)（即 \(\dot{f}_1,\dots ,\dot{f}_n\) 在 \(\hat{\mathbf{S}}(\mathrm{E})\) 中算出的乘积）变为 \((:f_1\dots f_n:)\)。假定 X 是 Souslin 空间，且存在由类属于 E 且分离 X 的点的函数组成的可数族 \((f_{n})\)。这时 \(\phi\) 是从 \(\hat{\mathbf{S}}(\mathrm{E})\) 到 \(\mathrm{L}_{\mathbb{R}}^{2}(\mathrm{X},\mu)\) 上的同构。

